
\begin{register}{H}{dcsr}{0x7B0}
\label{regdesc:DCSR}
\regfield{xdebugver}{4}{28}{{4}}%
\regfield{reserved}{12}{16}{{0}}%
\regfield{ebreakm}{1}{15}{{0}}%
\regfield{reserved}{2}{13}{{0}}%
\regfield{ebreaku}{1}{12}{{0}}%
\regfield{stepie}{1}{11}{{0}}%
\regfield{stopcount}{1}{10}{{0}}%
\regfield{stoptime}{1}{9}{{0}}%
\regfield{cause}{3}{6}{{0}}%
\regfield{reserved}{1}{5}{{0}}%
\regfield{mprven}{1}{4}{{0}}%
\regfield{reserved}{1}{3}{{0}}%
\regfield{step}{1}{2}{{0}}%
\regfield{prv}{2}{0}{{0}}%
\reglabel{Reset}\regnewline%
\begin{regdesc}
\begin{reglist}
\label{fielddesc:DCSRXEDBGVER}\item [xdebugver] Represents the debug version. \\
4: External debug support exists\\
(RO)
\label{fielddesc:DCSREBREAKM}\item [ebreakm] When 1, ebreak instructions in Machine Mode enter debug mode. \\
(R/W)
\label{fielddesc:DCSREBREAKU}\item [ebreaku] When 1, ebreak instructions in User/Application Mode enter debug mode. \\
(R/W)
\label{fielddesc:DCSRSTEPIE}\item [stepie] Configures interrupt control in debug mode.
0: Interrupts are disabled during single-stepping \\
1: Interrupts are enabled during single-stepping \\
(R/W) \\
Debugger must not change the value of this bit while the hart is running.
\label{fielddesc:DCSRSTOPCOUNT}\item [stopcount] Configures mcycle counter control in debug mode.
0: Increment counters as usual.\\
1: Not increment any counters while in debug mode or on ebreak instructions that cause entry into debug mode. These counters include the cycle and instret CSRs. \\
(R/W)
\label{fielddesc:DCSRSTOPTIME}\item [stoptime] Configures timer control in debug mode.\\
0: Increment timers as usual.\\
1: Not increment any hart-local timers while in debug mode. \\
(R/W)
\label{fielddesc:DCSRCAUSE}\item [cause] Explains why debug mode was entered. When there are multiple reasons to enter debug mode in a single cycle, the cause with the highest priority number is the one written.\\
1: An ebreak instruction was executed. (priority 3)\\
2: The Trigger Module caused a halt. (priority 4, highest)\\
3: haltreq was set. (priority 1)\\
4: The CPU core single-stepped because the step was set. (priority 0, lowest)\\
5: The hart halted directory out of reset due to resethaltreq. (priority 2)\\
Other values are reserved for future use.\\
(RO)
\label{fielddesc:DCSRMPRVEN}\item [mprven] Configures whether to enable the functionality of MPRV bit in mstatus CSR during debug mode.\\
0: Disable\\
1: Enable\\
(R/W)\\
\item [Continued on the next page...]
\end{reglist}\end{regdesc}
\end{register}
\addtocounter{Regfloat}{-1}
\begin{register}{H}{dcsr}{0x7B0}
\begin{regdesc}\begin{reglist}
\item [Continued from the previous page...]

\label{fielddesc:DCSRSTEP}\item [step] When set and not in debug mode, the core will only execute a single instruction and then enter debug mode. \\
If the instruction does not complete due to an exception, the core will immediately enter debug mode before executing the trap handler, with appropriate exception registers set.\\
Setting this bit does not mask interrupts. This is a deviation from the RISC-V External Debug Support Specification Version 0.13.\\
(R/W)\\
\label{fielddesc:DCSRPRV}\item [prv] Contains the privilege level the core was operating in when debug mode was entered. A debugger can change this value to change the core's privilege level when exiting debug mode. Only \textbf{0x3} (machine mode) and \textbf{0x0} (user mode) are supported. (R/W)
\end{reglist}
\end{regdesc}
\end{register}

\begin{register}{H}{dpc}{0x7B1}
\label{regdesc:DPC}
\regfield{dpc}{32}{0}{{0}}%
\reglabel{Reset}\regnewline%
\begin{regdesc}
\begin{reglist}
\label{fielddesc:DPC}\item [dpc] Upon entry to debug mode, dpc is written with the virtual address of the next instruction to be executed. When resuming, the CPU core's PC is updated to the virtual address stored in dpc. A debugger may write dpc to change where the CPU resumes. (R/W)
\end{reglist}
\end{regdesc}
\end{register}

\begin{longtable}{| l| p{12cm}|}
\caption{Virtual address in DPC upon Debug Mode Entry} \label{tab:riscvcpu-dbg-cause}
\\ \hline
\rowcolor{lightgray}
\textbf{Cause} & \textbf{Virtual Address in DPC}   \\ \hline
ebreak \label{ebreak} &  Address of ebreak instruction  \\ \hline
Single Step\label{sstep} & Address of the instruction that would be executed next if no debugging was going on i.e. pc+4 for 32-bit instruction that don't change program flow, the destination PC on taken jumps/branches, etc.  \\ \hline
Trigger Module \label{trigger} &  If \hyperref[fielddesc:MCONTROLTIMING]{timing} is 0, the address of the instruction which caused the trigger to fire. If \hyperref[fielddesc:MCONTROLTIMING]{timing} is 1, the address of next instruction to be executed at the time that debug mode was entered.\\ \hline
Halt Request\label{halt} & Address of the next instruction to be executed at the time that debug mode was entered.   \\ \hline
\end{longtable}

\begin{register}{H}{dscratch0}{0x7B2}
\label{regdesc:DSCRATCH0}
\regfield{dscratch0}{32}{0}{{0}}%
\reglabel{Reset}\regnewline%
\begin{regdesc}
\begin{reglist}
\label{fielddesc:DSCRATCH0}\item [dscratch0] Used by Debug Module internally. (R/W)
\end{reglist}
\end{regdesc}
\end{register}

\begin{register}{H}{dscratch1}{0x7B3}
\label{regdesc:DSCRATCH1}
\regfield{dscratch1}{32}{0}{{0}}%
\reglabel{Reset}\regnewline%
\begin{regdesc}
\begin{reglist}
\label{fielddesc:DSCRATCH1}\item [dscratch1] Used by Debug Module internally. (R/W)
\end{reglist}
\end{regdesc}
\end{register}